\documentclass[11pt,letterpaper]{article}
\usepackage[T1]{fontenc}
\usepackage{lmodern,microtype}
\usepackage[margin=0.78in,headheight=15pt,headsep=15pt]{geometry}
\usepackage{mathtools,amssymb}
\usepackage{xcolor,booktabs,array,tikz}
\usepackage[most]{tcolorbox}
\usepackage{enumitem,fancyhdr,needspace,hyperref}
\definecolor{ink}{HTML}{162332}
\definecolor{navy}{HTML}{17324D}
\definecolor{teal}{HTML}{167D82}
\definecolor{mist}{HTML}{EDF5F4}
\definecolor{muted}{HTML}{5D6B78}
\color{ink}
\hypersetup{colorlinks=true,linkcolor=teal,urlcolor=teal,pdfauthor={Igor Rivin}}
\setlength{\parindent}{0pt}
\setlength{\parskip}{4pt}
\setlength{\abovedisplayskip}{7pt}
\setlength{\belowdisplayskip}{7pt}
\setlist[enumerate]{label=(\alph*),leftmargin=1.8em,itemsep=3pt,topsep=4pt}
\newcommand{\Z}{\mathbb Z}
\newcommand{\F}{\mathbb F}
\newcommand{\one}{\mathbf{1}}
\newcommand{\id}{\operatorname{id}}
\newcommand{\pts}[1]{\hfill\textcolor{muted}{\small[#1 points]}}
\newcommand{\partpts}[1]{\unskip\nobreak\hspace{.4em}\textcolor{muted}{\small(#1)}}
\newcommand{\edition}{\ifsolutions Instructor solutions\else Student handout\fi}
\newcounter{problem}
\newcommand{\problem}[2]{\par\addvspace{11pt}\Needspace{6\baselineskip}%
  \stepcounter{problem}{\bfseries\color{navy}\theproblem.\ #1}\pts{#2}\par\nobreak}
\newcommand{\topic}[1]{\par\addvspace{9pt}\Needspace{10\baselineskip}%
  {\large\bfseries\color{teal}#1}\par\nobreak}
\newcommand{\studentbreak}{\ifsolutions\else\newpage\fi}
\newcommand{\challenge}[1]{\par\addvspace{12pt}\Needspace{7\baselineskip}%
  {\bfseries\color{teal}Optional challenge: #1}\pts{10 extra}\par\nobreak}
\pagestyle{fancy}
\fancyhf{}
\fancyhead[L]{\small\color{muted}\setname}
\fancyhead[R]{\small\color{muted}Patterns, Proofs, and Primes}
\fancyfoot[L]{\footnotesize\color{muted}\edition}
\fancyfoot[R]{\small\color{teal}\thepage}
\renewcommand{\headrulewidth}{0.3pt}
\NewDocumentEnvironment{solution}{+b}{\ifsolutions
  \begin{tcolorbox}[enhanced,breakable,colback=mist!45,colframe=teal!45,
    boxrule=.5pt,arc=1mm,left=3mm,right=3mm,top=2mm,bottom=2mm,
    before skip=7pt,after skip=7pt,title={Solution},
    title after break={Solution (continued)},colbacktitle=mist,
    coltitle=teal,fonttitle=\small\bfseries]
  \small #1
  \end{tcolorbox}\fi}{}
\newcommand{\handouttitle}[3]{%
  {\LARGE\bfseries\color{navy}\setname}\hfill{\small\color{teal}\edition}\\[4pt]
  {\large #1}\\[4pt]{\small\color{muted}#2}\par
  \medskip\textbf{Preparation.} #3\par
  \ifsolutions\else\medskip
  \textbf{Name:}\ \rule{.46\linewidth}{.4pt}\hfill
  \textbf{Due:}\ \rule{.20\linewidth}{.4pt}\par\fi}

\newcommand{\setname}{Multiplicative functions}
\newcommand{\instructorbreak}{\ifsolutions\newpage\fi}
\newcommand{\credit}[1]{\par\smallskip
  {\footnotesize\color{muted}\textit{Grading:} #1}}
\ifsolutions
  \setlength{\parskip}{3pt}
  \setlength{\abovedisplayskip}{5pt}
  \setlength{\belowdisplayskip}{5pt}
  \setlist[enumerate]{label=(\alph*),leftmargin=1.8em,
    itemsep=2pt,topsep=3pt,parsep=0pt}
\fi
\hypersetup{pdftitle={Number Theory: Multiplicative Functions Problem Set},
  pdfsubject={Prime powers, divisor functions, totients, convolution, and Mobius cancellation}}

\begin{document}
\handouttitle{A problem set on prime powers and divisor sums}
{8 problems \textbar\ 100 points + 10 optional challenge points \textbar\ Suggested time: 2--3 hours}
{Unique factorization, gcds, finite sums, and inclusion--exclusion.}

\textbf{Instructions.} Write on separate paper. Show factorizations and
justify your answers. When asked for a proof, do not merely cite the
statement. You may use results from earlier problems even if you have
not proved them. Computation may check examples but does not replace a proof.

\textbf{Conventions.} All integers and divisors in this assignment are
positive unless stated otherwise. An arithmetic function is a function
on the positive integers with real values. It is \emph{multiplicative} if
$f(1)=1$ and $f(mn)=f(m)f(n)$ whenever $\gcd(m,n)=1$; it is
\emph{completely multiplicative} if $f(1)=1$ and the product rule holds
for all $m,n$. Thus $f(\prod_i p_i^{a_i})=\prod_i f(p_i^{a_i})$
for multiplicative $f$ and distinct primes $p_i$. Empty products have value $1$.

\ifsolutions
\textbf{Grading note.} Subpart scores total 100; the challenge is extra.
Accept equivalent proofs and give method credit after a carried arithmetic
error. Check uses of coprimality and the value at $1$. The time estimate
excludes the challenge and assumes practice writing short proofs.
\fi

\topic{I. What prime powers tell us}

\problem{Read the definition carefully}{10}
Suppose $f$ is multiplicative and
\[
 f(2)=2,\qquad f(4)=5,\qquad f(3)=-1,\qquad
 f(9)=4,\qquad f(5)=0.
\]
\begin{enumerate}
  \item Find $f(1),f(12),f(18),f(36)$, and $f(45)$. \partpts{5}
  \item Do these data determine $f(8)$? Explain by describing two
  multiplicative functions satisfying all the given data but having
  different values at $8$. \partpts{3}
  \item Could $f$ be completely multiplicative? Justify your answer.
  \partpts{2}
\end{enumerate}
\begin{solution}
\begin{enumerate}
  \item The values, in order, are
  \[
    1,\quad f(4)f(3)=-5,\quad f(2)f(9)=8,\quad
    f(4)f(9)=20,\quad f(9)f(5)=0.
  \]
  The factors in each displayed product are coprime.
  \item No. Assign the five specified prime-power values as given.
  For a first function assign value $0$ at $2^3$; for a second assign
  value $7$ there. Assign value $1$ at every other unspecified prime
  power. Set each function's value at $1$ to $1$, and at a general
  integer take the product over its prime-power factors. Unique
  factorization makes both functions well defined and multiplicative,
  since coprime inputs have disjoint prime factors.
  \item No: complete multiplicativity would require
  $f(4)=f(2)^2=4$, but the given value is $5$.
\end{enumerate}
\credit{(a) One point per value. (b) Nonuniqueness: 1; two consistent
prime-power assignments: 1; multiplicative extension: 1.
(c) Required equality and contradiction: 2.}
\end{solution}

\instructorbreak
\problem{Which operations preserve multiplicativity?}{10}
Let $f$ and $g$ be multiplicative functions.
\begin{enumerate}
  \item Prove that the pointwise product $h(n)=f(n)g(n)$ is
  multiplicative. Prove the same for $q(n)=f(n)/g(n)$ if $g(n)\ne0$
  for every $n$. \partpts{4}
  \item Explain why the pointwise sum $f+g$ is never multiplicative
  under our convention. \partpts{2}
  \item Normalizing the sum fixes its value at $1$. Does it follow that
  $A(n)=(f(n)+g(n))/2$ is always multiplicative? Prove this or give
  an explicit counterexample with coprime inputs. \partpts{4}
\end{enumerate}
\begin{solution}
\begin{enumerate}
  \item Both $h(1)$ and $q(1)$ are $1$. For coprime $m,n$,
  \[
    h(mn)=f(m)f(n)g(m)g(n)=h(m)h(n),
    \qquad
    q(mn)=\frac{f(m)f(n)}{g(m)g(n)}=q(m)q(n).
  \]
  Nonvanishing of $g$ makes every quotient defined.
  \item $(f+g)(1)=2$, whereas a multiplicative function has value $1$
  at $1$.
  \item Take $f(n)=1$ and $g(n)=n$, which are both completely
  multiplicative. Then $A(n)=(n+1)/2$. Although $A(1)=1$,
  \[
    A(6)=\frac72\ne3=\frac32\cdot2=A(2)A(3),
    \qquad\gcd(2,3)=1.
  \]
\end{enumerate}
\credit{(a) Each proof, including its value at $1$: 2.
(b) Value at $1$ and conclusion: 2.
(c) Valid multiplicative examples: 1; normalized values: 1;
coprime inputs: 1; failed equality: 1.}
\end{solution}

\problem{Counting and adding divisors}{14}
Define $\tau(n)=\sum_{d\mid n}1$ and $\sigma(n)=\sum_{d\mid n}d$.
Write $n=\prod_{i=1}^r p_i^{a_i}$, with distinct primes $p_i$.
\begin{enumerate}
  \item Describe the divisors of $n$ using exponents. Deduce product
  formulas for $\tau(n)$ and $\sigma(n)$, and explain why both
  functions are multiplicative. \partpts{6}
  \item Compute $\tau(540)$ and $\sigma(540)$. \partpts{4}
  \item Prove that $\tau(n)$ is odd if and only if $n$ is a perfect
  square. Include $n=1$. \partpts{4}
\end{enumerate}
\begin{solution}
\begin{enumerate}
  \item Every divisor is uniquely $\prod_i p_i^{b_i}$ with
  $0\le b_i\le a_i$. Counting choices, or distributing a product of
  finite sums, gives
  \[
    \tau(n)=\prod_i(a_i+1),\qquad
    \sigma(n)=\prod_i(1+p_i+\cdots+p_i^{a_i}).
  \]
  At $1$ both products are empty and have value $1$. Coprime inputs
  have disjoint prime factors, so the formulas give multiplicativity.
  \item $540=2^2\cdot3^3\cdot5$, hence
  \[
    \tau(540)=3\cdot4\cdot2=24,\qquad
    \sigma(540)=(1+2+4)(1+3+9+27)(1+5)=1680.
  \]
  \item A product of integers is odd exactly when all factors are odd.
  Thus $\tau(n)$ is odd exactly when every $a_i+1$ is odd, equivalently
  when every $a_i$ is even. Unique factorization says this is exactly
  the condition for $n$ to be a square. Also $\tau(1)=1$ and $1=1^2$.
  Alternatively, pair divisors $d,n/d$; an unpaired divisor occurs
  exactly when $d^2=n$.
\end{enumerate}
\credit{(a) Unique exponent description: 2; each product formula: 1;
multiplicativity including $1$: 2. (b) Each calculation: 2.
(c) Parity criterion: 2; square criterion and both directions: 1;
case $1$: 1.}
\end{solution}

\studentbreak
\instructorbreak
\topic{II. Totients and divisor sums}

\problem{The totient from excluded multiples}{14}
Let $\varphi(n)$ count integers $1\le k\le n$ with $\gcd(k,n)=1$;
in particular, $\varphi(1)=1$.
\begin{enumerate}
  \item For a prime $p$ and $a\ge1$, prove
  $\varphi(p^a)=p^a-p^{a-1}$. \partpts{3}
  \item Apply inclusion--exclusion to the multiples of the distinct
  primes dividing $n$ to prove
  \[
    \varphi(n)=n\prod_{p\mid n}\left(1-\frac1p\right).
  \]
  Use this formula to prove that $\varphi$ is multiplicative.
  \partpts{6}
  \item Compute $\varphi(540)$. \partpts{2}
  \item Prove that $\varphi(n)$ is even for every $n>2$.
  \emph{Hint:} pair $k$ with $n-k$. \partpts{3}
\end{enumerate}
\begin{solution}
\begin{enumerate}
  \item An integer is not coprime to $p^a$ exactly when $p$ divides it.
  Among $1,\ldots,p^a$, there are $p^{a-1}$ such multiples. Subtract
  these from the $p^a$ integers.
  \item If the distinct prime divisors are $p_1,\ldots,p_r$, being
  coprime means avoiding every multiple of these primes. For a subset
  $S$ of the primes, the number divisible by all of them is
  $n/\prod_{p\in S}p$; the denominator divides $n$.
  Inclusion--exclusion therefore gives
  \[
    \varphi(n)=\sum_{S\subseteq\{p_1,\ldots,p_r\}}
      (-1)^{|S|}\frac{n}{\prod_{p\in S}p}
      =n\prod_{i=1}^r\left(1-\frac1{p_i}\right).
  \]
  For coprime $m,n$ the prime sets are disjoint; the product for $mn$
  separates into the products for $m$ and $n$, with leading factor
  $mn=m\cdot n$. Also $\varphi(1)=1$.
  \item $\varphi(540)=540(1-1/2)(1-1/3)(1-1/5)=144$.
  \item For $n>2$, every counted $k$ lies in $1,\ldots,n-1$.
  Since $\gcd(n-k,n)=\gcd(k,n)$, the map $k\mapsto n-k$
  pairs counted integers and is its own inverse. A fixed point would
  have $2k=n$, so $\gcd(k,n)=k$; coprimality would force $k=1,n=2$.
  This is excluded. All pairs have two distinct elements, so the count
  is even.
\end{enumerate}
\credit{(a) Exclude exactly the multiples of $p$: 1; count: 1;
subtraction: 1. (b) Intersection counts: 2; inclusion--exclusion and
product: 2; multiplicativity including $1$: 2. (c) Setup and value: 2.
(d) Pair preserves coprimality: 1; no fixed points for $n>2$: 1;
even count: 1.}
\end{solution}

\instructorbreak
\problem{Multiplication by a divisor sum}{14}
The \emph{Dirichlet convolution} of arithmetic functions is
\[
  (f*g)(n)=\sum_{d\mid n}f(d)g(n/d).
\]
Let $\one(n)=1$, $\id(n)=n$, and let $\varepsilon(1)=1$,
$\varepsilon(n)=0$ for $n>1$.
\begin{enumerate}
  \item Suppose $\gcd(m,n)=1$. Prove that every divisor of $mn$
  has a unique expression $ab$ with $a\mid m$ and $b\mid n$.
  \partpts{4}
  \item Use (a) to prove that $f*g$ is multiplicative when $f$ and
  $g$ are multiplicative. Check the value at $1$ and explain each use
  of coprimality. \partpts{6}
  \item Identify $\one*\one$ and $\one*\id$ in terms of the
  functions already studied. For an arbitrary arithmetic function $f$,
  prove $f*\varepsilon=f$. \partpts{4}
\end{enumerate}
\begin{solution}
\begin{enumerate}
  \item No prime divides both $m$ and $n$. Assign each prime power in
  a divisor of $mn$ to the factor coming from $m$ or from $n$.
  This gives the required $a,b$ and makes them unique. Conversely,
  every such product $ab$ divides $mn$.
  \item Set $h=f*g$. We have $h(1)=f(1)g(1)=1$. For coprime $m,n$,
  both $\gcd(a,b)=1$ and $\gcd(m/a,n/b)=1$, so
  \[
  \begin{aligned}
    h(mn)&=\sum_{a\mid m}\sum_{b\mid n}f(ab)g((m/a)(n/b))\\
      &=\sum_{a\mid m}\sum_{b\mid n}f(a)f(b)g(m/a)g(n/b)\\
      &=\left(\sum_{a\mid m}f(a)g(m/a)\right)
        \left(\sum_{b\mid n}f(b)g(n/b)\right)=h(m)h(n).
  \end{aligned}
  \]
  \item $(\one*\one)(n)=\sum_{d\mid n}1=\tau(n)$.
  Also $(\one*\id)(n)=\sum_{d\mid n}n/d=\sigma(n)$, since
  $d\mapsto n/d$ permutes the divisors. Finally, only the term $d=n$
  survives in $\sum_{d\mid n}f(d)\varepsilon(n/d)$, giving $f(n)$.
\end{enumerate}
\credit{(a) Existence: 2; uniqueness and converse: 2.
(b) Value at $1$: 1; divisor reindexing: 1; both coprimality
conditions: 2; product factorization: 2.
(c) $\tau$: 1; $\sigma$: 1; identity proof: 2.}
\end{solution}

\instructorbreak
\problem{M\"obius cancellation}{14}
An integer is \emph{squarefree} if no prime square divides it. Define
\[
  \mu(n)=
  \begin{cases}
    1,&n=1,\\
    (-1)^r,&n\text{ is a product of }r\text{ distinct primes},\\
    0,&n\text{ is not squarefree}.
  \end{cases}
\]
\begin{enumerate}
  \item Prove that $\mu$ is multiplicative. Show that it is not
  completely multiplicative. \partpts{3}
  \item Find $\sum_{d\mid p^a}\mu(d)$ for a prime $p$ and $a\ge1$.
  Use Problem 5 to prove $\mu*\one=\varepsilon$, including its
  value at $1$. \partpts{5}
  \item You may use associativity and commutativity of convolution.
  Prove the following identities for every $n\ge1$. \partpts{6}
  \[
    \sum_{d\mid n}\mu(d)\sigma(n/d)=n,
    \qquad
    \sum_{d\mid n}\mu(d)\tau(n/d)=1.
  \]
\end{enumerate}
\begin{solution}
\begin{enumerate}
  \item We have $\mu(1)=1$. If either coprime input has a prime square
  factor, both $\mu(mn)$ and $\mu(m)\mu(n)$ vanish. Otherwise their
  disjoint lists of primes concatenate, so the signs multiply.
  Complete multiplicativity fails since $\mu(p^2)=0$ but
  $\mu(p)^2=1$.
  \item The prime-power sum is $\mu(1)+\mu(p)=1-1=0$;
  all higher powers contribute zero. By Problem 5,
  $H=\mu*\one$ is multiplicative. Thus for $n>1$, its product
  over prime-power factors contains a zero, so $H(n)=0$.
  At $1$, $H(1)=1$. This is precisely $\varepsilon$.
  \item Problem 5 gives $\sigma=\one*\id$ and
  $\tau=\one*\one$. Hence
  \[
    \mu*\sigma=(\mu*\one)*\id=\varepsilon*\id=\id,
    \qquad
    \mu*\tau=(\mu*\one)*\one=\varepsilon*\one=\one.
  \]
  Evaluating these function identities at $n$ gives the two requested
  sums. They also hold at $1$.
\end{enumerate}
\credit{(a) Multiplicativity with both cases and $1$: 2;
counterexample to complete multiplicativity: 1.
(b) Prime-power sum: 2; multiplicativity and reduction to prime
powers: 2; value at $1$: 1.
(c) Each convolution calculation and its translation to the sum: 3.}
\end{solution}

\studentbreak
\instructorbreak
\topic{III. Two applications}

\problem{An alternating divisor sum detects squares}{12}
For $n=\prod_{i=1}^r p_i^{a_i}$ define
$\Omega(n)=a_1+\cdots+a_r$, with $\Omega(1)=0$, and set
$\lambda(n)=(-1)^{\Omega(n)}$. Thus $\Omega$ counts prime factors
\emph{with multiplicity}.
\begin{enumerate}
  \item Prove that $\lambda$ is completely multiplicative.
  \partpts{4}
  \item Evaluate $\sum_{j=0}^a\lambda(p^j)$ and deduce, for every $n$,
  \[
    \sum_{d\mid n}\lambda(d)=
    \begin{cases}
      1,&n\text{ is a perfect square},\\
      0,&n\text{ is not a perfect square}.
    \end{cases}
  \]
  Justify the passage from prime powers to general $n$, and include
  $n=1$. \partpts{6}
  \item Evaluate $\sum_{d\mid 360}\lambda(d)$ and
  $\sum_{d\mid 900}\lambda(d)$ without listing all the divisors.
  \partpts{2}
\end{enumerate}
\begin{solution}
\begin{enumerate}
  \item Prime exponents add in a product whether or not its factors
  are coprime, so $\Omega(mn)=\Omega(m)+\Omega(n)$. Therefore
  $\lambda(mn)=(-1)^{\Omega(m)+\Omega(n)}=\lambda(m)\lambda(n)$
  for all $m,n$, and $\lambda(1)=1$.
  \item The prime-power sum is $1-1+\cdots+(-1)^a$, equal to
  $1$ for even $a$ and $0$ for odd $a$. By Problem 5 the divisor-sum
  function $\lambda*\one$ is multiplicative. Its value at
  $\prod_i p_i^{a_i}$ is the product of these alternating sums,
  which is $1$ exactly when every $a_i$ is even, and $0$ otherwise.
  This is the square criterion from unique factorization.
  At $1$, the divisor sum is $\lambda(1)=1$.
  \item $360=2^3\cdot3^2\cdot5$ is not a square, so its sum is $0$.
  Since $900=2^2\cdot3^2\cdot5^2=30^2$, its sum is $1$.
\end{enumerate}
\credit{(a) Exponents add without a coprimality restriction: 2;
sign product and value at $1$: 2. (b) Alternating sum: 2;
multiplicativity and product: 2; square criterion: 1; case $1$: 1.
(c) Each evaluated sum: 1.}
\end{solution}

\instructorbreak
\problem{A sum of gcds}{12}
Define $G(n)=\sum_{k=1}^n\gcd(k,n)$.
\begin{enumerate}
  \item Group the integers $k$ according to $d=\gcd(k,n)$ to prove
  \[
    G(n)=\sum_{d\mid n}d\,\varphi(n/d).
  \]
  Explain why the group with gcd equal to $d$ has $\varphi(n/d)$
  members. \partpts{5}
  \item Use the earlier problems to deduce that $G$ is multiplicative.
  \partpts{2}
  \item For a prime $p$ and $a\ge1$, derive
  \[
    G(p^a)=(a+1)p^a-a p^{a-1}.
  \]
  Then compute $G(60)$ without summing 60 gcds individually.
  \partpts{5}
\end{enumerate}
\begin{solution}
\begin{enumerate}
  \item The group with gcd $d$ consists exactly of $k=dj$ with
  $1\le j\le n/d$ and $\gcd(j,n/d)=1$. There are
  $\varphi(n/d)$ choices, each contributing $d$ to the sum.
  These groups partition $1,\ldots,n$, so their contributions add
  to the formula. The case $d=n$ contains only $k=n$, consistent
  with $\varphi(1)=1$.
  \item The formula says $G=\id*\varphi$; both factors are
  multiplicative, so Problem 5 applies. In particular $G(1)=1$.
  \item The term $d=p^a$ contributes $p^a\varphi(1)=p^a$.
  For $d=p^j$ with $0\le j<a$, the contribution is
  \[
    p^j\varphi(p^{a-j})
      =p^j(p^{a-j}-p^{a-j-1})=p^a-p^{a-1}.
  \]
  There are $a$ such terms, giving
  $G(p^a)=p^a+a(p^a-p^{a-1})$ as required. Finally,
  \[
    G(60)=G(4)G(3)G(5)
      =(3\cdot4-2\cdot2)(2\cdot3-1)(2\cdot5-1)
      =8\cdot5\cdot9=360.
  \]
\end{enumerate}
\credit{(a) Gcd class bijection: 3; weighted contribution: 1;
partition and sum: 1. (b) Convolution and multiplicativity: 2.
(c) Prime-power derivation including its final term: 3;
factorization and value at $60$: 2.}
\end{solution}

\challenge{Even perfect numbers}
An integer $N$ is \emph{perfect} if its positive divisors other than $N$
sum to $N$, equivalently $\sigma(N)=2N$.
\begin{enumerate}
  \item Let $r\ge2$ and suppose $2^r-1$ is prime. Prove that
  $N=2^{r-1}(2^r-1)$ is perfect. Give the examples from $r=2,3,5$.
  \partpts{4}
  \item Prove that every even perfect number has the form in (a).
  \emph{Hint:} write $N=2^a m$ with $a\ge1$ and $m$ odd, and put
  $q=2^{a+1}-1$. First show $q\mid m$. Write $m=qt$ and show
  $\sigma(m)=m+t$. If $t>1$, compare this with the contribution of
  the three distinct divisors $1,t,m$. \partpts{6}
\end{enumerate}
\begin{solution}
\begin{enumerate}
  \item Set $q=2^r-1$, an odd prime. The factors $2^{r-1}$ and
  $q$ are coprime, so Problem 3 gives
  \[
    \sigma(N)=(1+2+\cdots+2^{r-1})(1+q)
      =(2^r-1)2^r=2N.
  \]
  The examples are $6,28,496$.
  \item The perfect-number equation is
  \[
    (2^{a+1}-1)\sigma(m)=2^{a+1}m,
    \quad\text{or}\quad q\sigma(m)=(q+1)m.
  \]
  Since $\gcd(q,q+1)=1$, Euclid's lemma gives $q\mid m$.
  Write $m=qt$. Dividing by $q$ gives
  $\sigma(m)=(q+1)t=m+t$. We have $q\ge3$, so $t<m$.
  If $t>1$, the three distinct divisors $1,t,m$ imply
  $\sigma(m)\ge1+t+m>m+t$, a contradiction. Thus $t=1$,
  $m=q>1$, and $\sigma(m)=m+1$. A composite $m$ would have
  an additional positive divisor besides $1,m$, so $m$ must be
  prime. Setting $r=a+1\ge2$ gives the required form.
\end{enumerate}
\credit{(a) Coprime factorization and divisor-sum formula: 2;
perfect-number equality: 1; examples: 1. (b) Equation and $q\mid m$:
2; substitution and divisor contradiction: 2; $t=1$, primality,
and the final form: 2.}
\end{solution}

\end{document}
